\documentclass[10pt,a4paper]{amsart}
\usepackage{amsmath,amssymb,amsthm,graphicx}
% Retained minimal preamble for the Stacks Project chapter more-algebra.tex (More on Algebra).
% The chapter file pinned-original-source.tex carries only \input{preamble} on its first line, so no
% preamble text lives inside the chapter itself.  The Stacks Project uses one single shared preamble for
% every chapter; the snapshot stages it identically as sources/stacks-chow/preamble.tex and
% sources/stacks-etale-cohomology/preamble.tex.  The block below is copied mechanically, line for line,
% from sources/stacks-chow/preamble.tex lines 14-18 (the Xy-pic setup) and 154-217 (the native theorem
% declarations together with the mathematical macros).  Nothing was added, removed or redefined; the
% wrapper class and the amsmath/amssymb/amsthm/graphicx base packages are supplied by the assembler.
% Chapter: more-algebra.tex -- More on Algebra.
% For commutative diagrams we use Xy-pic
\usepackage[all]{xy}

% We use 2cell for 2-commutative diagrams.
\xyoption{2cell}

% Theorem environments.
%
\theoremstyle{plain}
\newtheorem{theorem}[subsection]{Theorem}
\newtheorem{proposition}[subsection]{Proposition}
\newtheorem{lemma}[subsection]{Lemma}

\theoremstyle{definition}
\newtheorem{definition}[subsection]{Definition}
\newtheorem{example}[subsection]{Example}
\newtheorem{exercise}[subsection]{Exercise}
\newtheorem{situation}[subsection]{Situation}

\theoremstyle{remark}
\newtheorem{remark}[subsection]{Remark}
\newtheorem{remarks}[subsection]{Remarks}

\numberwithin{equation}{subsection}

% Macros
%
\def\lim{\mathop{\mathrm{lim}}\nolimits}
\def\colim{\mathop{\mathrm{colim}}\nolimits}
\def\Spec{\mathop{\mathrm{Spec}}}
\def\Hom{\mathop{\mathrm{Hom}}\nolimits}
\def\Ext{\mathop{\mathrm{Ext}}\nolimits}
\def\SheafHom{\mathop{\mathcal{H}\!\mathit{om}}\nolimits}
\def\SheafExt{\mathop{\mathcal{E}\!\mathit{xt}}\nolimits}
\def\Sch{\mathit{Sch}}
\def\Mor{\mathop{\mathrm{Mor}}\nolimits}
\def\Ob{\mathop{\mathrm{Ob}}\nolimits}
\def\Sh{\mathop{\mathit{Sh}}\nolimits}
\def\NL{\mathop{N\!L}\nolimits}
\def\CH{\mathop{\mathrm{CH}}\nolimits}
\def\proetale{{pro\text{-}\acute{e}tale}}
\def\etale{{\acute{e}tale}}
\def\QCoh{\mathit{QCoh}}
\def\Ker{\mathop{\mathrm{Ker}}}
\def\Im{\mathop{\mathrm{Im}}}
\def\Coker{\mathop{\mathrm{Coker}}}
\def\Coim{\mathop{\mathrm{Coim}}}

% Boxtimes
%
\DeclareMathSymbol{\boxtimes}{\mathbin}{AMSa}{"02}

%
% Macros for moduli stacks/spaces
%
\def\QCohstack{\mathcal{QC}\!\mathit{oh}}
\def\Cohstack{\mathcal{C}\!\mathit{oh}}
\def\Spacesstack{\mathcal{S}\!\mathit{paces}}
\def\Quotfunctor{\mathrm{Quot}}
\def\Hilbfunctor{\mathrm{Hilb}}
\def\Curvesstack{\mathcal{C}\!\mathit{urves}}
\def\Polarizedstack{\mathcal{P}\!\mathit{olarized}}
\def\Complexesstack{\mathcal{C}\!\mathit{omplexes}}
% \Pic is the operator that assigns to X its picard group, usage \Pic(X)
% \Picardstack_{X/B} denotes the Picard stack of X over B
% \Picardfunctor_{X/B} denotes the Picard functor of X over B
\def\Pic{\mathop{\mathrm{Pic}}\nolimits}
\def\Picardstack{\mathcal{P}\!\mathit{ic}}
\def\Picardfunctor{\mathrm{Pic}}
\def\Deformationcategory{\mathcal{D}\!\mathit{ef}}

% --- appended 2026-10-12 (Chief): the project own definitions of the reference/slogan/history
% display environments, copied mechanically from the shared Stacks preamble lines 8-12
% (\usepackage{verbatim} plus the three \newenvironment lines, which the chapter preamble keeps
% verbatim from the project preamble). Nothing added, removed or redefined: these environments
% carry no mathematical content and the project itself typesets them as comments. Without these
% three lines a counted statement whose own text uses them fails with Environment slogan undefined.
\usepackage{verbatim}
\newenvironment{reference}{\comment}{\endcomment}
\newenvironment{slogan}{\comment}{\endcomment}
\newenvironment{history}{\comment}{\endcomment}
\usepackage{cleveref}
\crefname{originallemma}{original lemma}{original lemmas}
\Crefname{originallemma}{Original lemma}{Original lemmas}
\makeatletter
\newlabel{lemma-neighbourhood-isomorphism}{{05E9}{source}{}{sourcepointer-1}{}}
\newlabel{lemma-neighbourhood-isomorphism@cref}{{[originallemma][0][]05E9}{[1][0][]source}{}{}{}}
\makeatother
\crefname{originallemma}{original lemma}{original lemmas}
\Crefname{originallemma}{Original lemma}{Original lemmas}
\makeatletter
\newlabel{lemma-homotopy-koszul}{{0663}{source}{}{sourcepointer-2}{}}
\newlabel{lemma-homotopy-koszul@cref}{{[originallemma][0][]0663}{[1][0][]source}{}{}{}}
\makeatother
\crefname{originallemma}{original lemma}{original lemmas}
\Crefname{originallemma}{Original lemma}{Original lemmas}
\makeatletter
\newlabel{lemma-extended-alternating-torsion}{{0G6K}{source}{}{sourcepointer-3}{}}
\newlabel{lemma-extended-alternating-torsion@cref}{{[originallemma][0][]0G6K}{[1][0][]source}{}{}{}}
\makeatother
\title{Koszul and alternating Cech complexes detect quasi-isomorphisms modulo I for maps identifying R/I and S/IS}
\author{The Stacks Project Authors}
\date{stacks-project:more-algebra.tex:local-snapshot-20261012; GNU FDL 1.2 or later}
\begin{document}
\maketitle
\setcounter{subsection}{1}
\noindent\textit{Editorial note.} All delivered mathematical text is the author's own, copied from the pinned Stacks Project chapter file more-algebra.tex; no statement or proof step was rewritten, reordered or removed. The counted span is the author environment and its own proof environment. References inside the span are handed to the assembler as navigation stubs that carry the canonical tag of the cited tag only; whether each of them is load-bearing is recorded on the stub itself.\par\medskip
\noindent\textit{Allowed standard prerequisites (editorial).}\begin{itemize}
\item A map of complexes K -> K tensor S over a flat ring map is a quasi-isomorphism as soon as the homology groups of K are I-power torsion and the map induces an isomorphism R/I -> S/IS on a neighbourhood. invoked at proof lines 24386-24388. Reduces both claims to the torsion property of the homology.
\item The homology groups of the Koszul complex K(R, f\_1, ..., f\_r) are I-power torsion when V(f\_1, ..., f\_r) = V(I). invoked at proof lines 24388-24389. Gives the torsion property for the Koszul case.
\item The homology groups of the extended alternating Cech complex on f\_1, ..., f\_r are I-power torsion. invoked at proof lines 24389-24391. Gives the torsion property for the Cech case.
\end{itemize}\medskip
\ifdefined\hypertarget\hypertarget{sourcepointer-1}{}\fi
\noindent\textit{Source reference.} Original lemma 05E9: sources/stacks-more-algebra/pinned-original-source.tex, source lines 24264--24264. Cited at proof lines 24386-24387 for the reduction to I-power torsion homology; the conclusion is used inside the counted span.\par\medskip
\ifdefined\hypertarget\hypertarget{sourcepointer-2}{}\fi
\noindent\textit{Source reference.} Original lemma 0663: sources/stacks-more-algebra/pinned-original-source.tex, source lines 6660--6660. Cited at proof lines 24388-24389 for the torsion property in the Koszul case; the conclusion is used inside the counted span.\par\medskip
\ifdefined\hypertarget\hypertarget{sourcepointer-3}{}\fi
\noindent\textit{Source reference.} Original lemma 0G6K: sources/stacks-more-algebra/pinned-original-source.tex, source lines 7066--7066. Cited at proof lines 24389-24391 for the torsion property in the Cech case; the conclusion is used inside the counted span.\par\medskip
\begin{lemma}
\label{lemma-map-identifies-koszul-and-cech-complexes}
Assume $\varphi : R \to S$ is a flat ring map and $I \subset R$ is a
finitely generated ideal such that $R/I \to S/IS$ is an isomorphism.
For any $f_1, \ldots, f_r \in R$ such that $V(f_1, \ldots, f_r) = V(I)$
\begin{enumerate}
\item the map of Koszul complexes
$K(R, f_1, \ldots, f_r) \to K(S, f_1, \ldots, f_r)$ is a quasi-isomorphism, and
\item The map of extended alternating {\v C}ech complexes
$$
\xymatrix{
R \to \prod_{i_0} R_{f_{i_0}} \to \prod_{i_0 < i_1} R_{f_{i_0}f_{i_1}}
\to \ldots \to R_{f_1\ldots f_r} \ar[d] \\
S \to \prod_{i_0} S_{f_{i_0}} \to \prod_{i_0 < i_1} S_{f_{i_0}f_{i_1}}
\to \ldots \to S_{f_1\ldots f_r}
}
$$
is a quasi-isomorphism.
\end{enumerate}
\end{lemma}

\begin{proof}
In both cases we have a complex $K_\bullet$ of $R$ modules and we want
to show that $K_\bullet \to K_\bullet \otimes_R S$ is a quasi-isomorphism.
By Lemma \ref{lemma-neighbourhood-isomorphism} and the flatness of
$R \to S$ this will hold as soon as all homology groups of $K$ are $I$-power
torsion. This is true for the Koszul complex by
Lemma \ref{lemma-homotopy-koszul} and for the extended alternating {\v C}ech
complex by Lemma \ref{lemma-extended-alternating-torsion}.
\end{proof}

\end{document}
